\section{最终结论}

面向后续实验，结论如下。

\begin{enumerate}
    \item 在现实参考范围内，静态 UEP/ICCP 场幅值最大，但识别特征弱。它更适合用来验证准静态场、空间梯度、多通道一致性和慢变包络，不是当前最容易做出清晰频谱特征的目标。
    \item 轴频场幅值约为静态场的 $m_1$ 倍。$m_1=0.03$ 时，轴频基频比静态场低 $30.5\ \mathrm{dB}$，但由于它有稳定线谱，是最适合当前实验验证的目标。
    \item 100 m 等级适合做链路验证。该距离下信号强，适合验证 3 Hz、6 Hz、9 Hz 注入、采集同步、频谱检测、积分增益和通道一致性。
    \item 300 m 等级开始考验噪声底。若系统等效噪声谱密度偏高，轴频基频会从稳健可测进入临界可测，需要窄带积分和更稳定的低频通道。
    \item 1000 m 等级必须依赖低噪声、长积分和稳定通道。静态场仍可能可测，但轴频基频通常是更严格的系统能力测试。
    \item 尾流不能与远场偶极子直接排序。当前只保留 $E_{\mathrm{wake,local}}\le vB_0$ 的局部上限，后续若要研究尾流，应单独建立流体-电磁耦合模型。
    \item 工频泄漏只作为实验干扰，不作为真实舰艇主源项。实验中必须记录 $A_{50,i}$、相位和谐波，并通过屏蔽、接地、差分和窄带剔除处理。
\end{enumerate}

因此，下一步实验建议按以下顺序推进：先做 3 Hz、6 Hz、9 Hz 可控低频电偶极子注入，目标场强覆盖 $10^{-6}$ 到 $10^{-8}\ \mathrm{V/m}$；再做准静态/慢变等效场；同步测量和剔除 50 Hz 干扰；尾流留作后续独立建模问题。
